\section{ Grundlagen der Vektor- und Matrizenrechnung}

\begin{ejercicio}[Vektoroperationen]
    Gegeben die Vektoren $\vec{a} = \begin{pmatrix}
        2\\
        5\\
        3
    \end{pmatrix}$, $\vec{b} =~\begin{pmatrix}
        -1\\
        5\\
        8
    \end{pmatrix}$ und $\vec{c} = \begin{pmatrix}
        5\\
        0\\
        1
    \end{pmatrix}$.
    Berechnen Sie die Komponenten (Vektorkoordinaten) sowie Betrag und Richtungswinkel des Vektors
    \begin{equation*}
        \vec{s} = -3\vec{b} + 4 (\vec{a} \cdot \vec{c}) \vec{a} - (\vec{a} -\vec{b}) \cdot (2 \vec{c} + \vec{a}) \vec{c}
    \end{equation*}
\end{ejercicio}

\begin{ejercicio}[Produkte mit Vektoren]
    Gegeben die Vektoren $\vec{a} = \begin{pmatrix}
        2\\
        1\\
        5
    \end{pmatrix}$, $\vec{b} =~\begin{pmatrix}
        -1\\
        3\\
        1
    \end{pmatrix}$ und $\vec{c} = \begin{pmatrix}
        4\\
        5\\
        1
    \end{pmatrix}$.
    Berechnen Sie mit diesen Vektoren die folgenden Produkte:
    \begin{enumerate}
        \item $\vec{a} \cdot \vec{b}$
        \item $(\vec{a} - \vec{c}) \cdot (\vec{a} + \vec{c})$
        \item $\vec{b} \times \vec{c}$
        \item $(\vec{a} \times \vec{c}) \cdot \vec{b}$
    \end{enumerate}
\end{ejercicio}

\begin{ejercicio}[Lineare (Un-)Abhängigkeit]
    Stellen Sie fest, ob die Vektoren $\vec{a}$, $\vec{b}$ und $\vec{c}$ linear unabhängig oder linear abhängig sind:
    \begin{enumerate}
        \item $\vec{a} = \begin{pmatrix}
            2\\
            1\\
            5
        \end{pmatrix}$, $\vec{b} = \begin{pmatrix}
            1\\
            4\\
            1
        \end{pmatrix}$ und $\vec{c} = \begin{pmatrix}
            -2\\
            4\\
            1
        \end{pmatrix}$
        \item $\vec{a} = \begin{pmatrix}
            5\\
            1\\
            1
        \end{pmatrix}$, $\vec{b} = \begin{pmatrix}
            -1\\
            0\\
            5
        \end{pmatrix}$ und $\vec{c} = \begin{pmatrix}
            -13\\
            -2\\
            13
        \end{pmatrix}$
    \end{enumerate}
\end{ejercicio}

\begin{ejercicio}[Orthogonalität von Vektoren]
    Eine Kraft $\vec{F} = \begin{pmatrix}
        F_x\\
        F_y\\
        F_z
    \end{pmatrix}$ vom Betrag $F = 100\sqrt{14} \ \mathrm{N}$ soll senkrecht auf einer Ebene $E$ stehen, die durch die Vektoren $\vec{a} = \begin{pmatrix}
        1\\
        -1\\
        1
    \end{pmatrix}$ und $\vec{b} = \begin{pmatrix}
        2\\
        1\\
        1
    \end{pmatrix}$ aufgespannt wird.
    Bestimmen Sie die (skalaren) Komponenten $F_x$, $F_y$ und $F_z$ dieser Kraft. Wieviele Lösungen gibt es?
\end{ejercicio}

\begin{ejercicio}[Matrizenmultiplikation]~
    \begin{enumerate}
        \item Zeigen Sie am Beispiel der folgenden beiden Matrizen
        \begin{equation*}
            A = \begin{pmatrix}
                -6 & 9 & -3\\
                -4 & 6 & -2\\
                -2 & 3 & -1
            \end{pmatrix} \quad \text{und} \quad B = \begin{pmatrix}
                0 & 1 & -2\\
                -1 & 0 & 3\\
                2 & -3 & 0
            \end{pmatrix},
        \end{equation*}
        dass die Matrizenmultiplikaton nicht-kommutativ ist, d.h. $A \cdot B \neq B \cdot A$.
        \item Gegeben die folgenden 3 Matrizen:
        \begin{equation*}
            A_{(3,2)} = \begin{pmatrix}
                1 & 3\\
                2 & 0\\
                1 & 4
            \end{pmatrix},\qquad B_{(2,2)} = \begin{pmatrix}
                2 & 1\\
                3 & -2
            \end{pmatrix} \quad \text{und} \quad C_{(2,3)} = \begin{pmatrix}
                2 & 0 & -1\\
                6 & 1 & 3
            \end{pmatrix}.
        \end{equation*}
        \begin{enumerate}
            \item Bilden Sie alle möglichen Produkte $X\cdot Y$ aus zwei Faktoren.
            \item Welche Produkte $X\cdot Y\cdot Z$ aus drei verschiedenen Faktoren können Sie damit bilden?
        \end{enumerate}
    \end{enumerate}
\end{ejercicio}

\begin{ejercicio}[Lösungsmengen linearer Gleichungssysteme]
    Bestimmen Sie die Lösungsmenge des Gleichungssystems $A\vec{x} = \vec{b}$ mit
    \begin{equation*}
        A = \begin{pmatrix}
            1 & -1 & 2 & 1 & 0\\
            -2 & 4 & -2 & -1 & 3\\
            0 & -2 & -2 & -1 & -3\\
            1 & -2 & 2 & 1 & 1
        \end{pmatrix}, \quad \vec{b} = \begin{pmatrix}
            1\\
            4\\
            -6\\
            1
        \end{pmatrix}.
    \end{equation*}
    Was können Sie zur Lösbarkeit des Gleichungssystems sagen?
\end{ejercicio}